% Este archivo es parte de la memoria del proyecto fin de carrera
% de Manuel López Urbina. Protegida bajo la licencia GFDL.
% Para más información, la licencia completa viene incluida en el
% fichero fdl-1.3.tex

% Copyright (C) 2012 Manuel López Urbina

\chapter{Introducción}
\label{chap:introducción}

\section{Introducción y antecedentes}

En la actualidad, los robots comerciales e industriales son ampliamente utilizados y realizan tareas de forma más exacta o más barata que los humanos. También se les utiliza en trabajos demasiado sucios, peligrosos o tediosos. Los robots son muy utilizados en plantas de fabricación, montaje y embalaje, en transporte, en exploraciones en la Tierra y en el espacio, cirugía, armamento, investigación en laboratorios y en la producción en masa de bienes industriales o de consumo. Otras aplicaciones incluyen la limpieza de residuos tóxicos, minería, búsqueda y rescate de personas y localización de minas terrestres, en definitiva, la robótica está presente en prácticamente cualquier ámbito de la actualidad.\\

Por otra parte, ninguno de los sistemas robóticos actuales podrían ser funcionales sin un software adecuado para automatización de tareas, siendo esencial una correcta sincronización entre los diferentes elementos hardware y software implicados con la finalidad de garantizar un correcto funcionamiento del conjunto robótico.\\

Centrándonos en lo anterior y aplicándolo en el campo de la automoción, ocasiona que en la actualidad se empleen importantes esfuerzos por parte de grupos investigadores y fabricantes de automóviles para dotar sus vehículos de los elementos necesarios para proporcionar a sus vehículos de un sistema de conducción autónoma. Campo en la actualidad en pleno desarrollo y presentando multitud problemas aún por resolver.\\

Entre los problemas existentes en los sistemas de conducción autónoma destaca su escasa adaptabilidad en carreteras reales donde existen multitud de imperfecciones, como pueden ser el desgaste o la falta de las señales viales, falta de iluminación o existencia de señalización provisional no fija como las indicadoras de obras en la vía. Estas circunstancias especiales impiden el correcto funcionamiento de los vehículos de dentro de unos niveles de seguridad mínimos aceptables para su implantación en carreteras transitables junto con vehículos conducidos de modo tradicional.\\

La problemática actual presentada, junto con que robótica y automoción son mis dos pasiones, hizo que me lanzara a la elaboración de este proyecto que unifica ambos campos anteriormente citados.\\

Así surgió \emph{OpenTSR}.\\

Se trata de un vehículo controlado vía WiFi dotado de una cámara de vídeo la cual visualiza su en torno en busca de aquellas señales de tráfico presentes y actuar en consecuencia.\\

\begin{figure}[H]
  \begin{center}
    \includegraphics[scale=0.3]{lente1.png}
  \end{center}
  \label{fig:logo}
 \caption{Logo OpenTSR \protect\footnotemark.}
\end{figure}

\footnotetext{Logotipo elaborado por Noelia Sales Montes para el proyecto OpenTSR.}

\emph{OpenTSR} (que será el nombre del sistema resultante) será una combinación de un elemento hardware (vehículo entre otros) y software (consola de control) surgido como una aproximación a la solución de los citados problemas.\\

\section{Objetivos}
\label{sec:objetivos}

El objetivo principal del proyecto es elaborar un software de detección de señales de tráfico mediante el análisis de las imágenes obtenidas a partir de una cámara. Dicho software irá asociado a un vehículo dotado de una cámara para permitir la visualización del terreno además de ser capaz de reconocer las señales de tráfico mostrando en el ordenador la última señal reconocida. Por otra parte, permitirá al vehículo actuar en consecuencia a la señal detectada. En definitiva se desea dotar al vehículo de un sistema de conducción autónoma \footnote{En la sección \ref{sec:def-vehículo-autónomo} referente a conceptos básicos se incluye la definición de vehículo autónomo.} a partir de señales de tráfico.\\

El sistema a desarrollar, por tanto, dispondrá de dos modos de funcionamiento, el primero de ellos proporciona una conducción a gusto del usuario, el cual podrá controlar a su antojo el vehículo por el entorno que le rodea mediante un joystick o haciendo uso de las teclas del teclado de un ordenador. En el segundo modo de funcionamiento, el vehículo actuará automáticamente según las señales de tráfico detectadas por la cámara, siendo capaz de ajustar velocidad o realizar giros y paradas de manera automática según los elementos visionados a través de la cámara.\\

\section{Acerca de este documento}

El documento se ha sido elaborado en un lenguaje sencillo y claro para permitir que un estudiante universitario de cualquier Ingeniería Informática pueda comprender los contenidos sin apenas dificultad añadida.\\

Este documento se organiza en los siguientes capítulos:\\

\begin{itemize}

\item En el capítulo \ref{chap:introducción}, Introducción, se comentan las razones que han motivado la creación de este proyecto, así como el propósito del mismo.

\item En el capítulo \ref{chap:conceptos-básicos}, Conceptos básicos, se incluyen definiciones de aquellos conceptos considerados de interés para la correcta comprensión del contenido de la presente memoria.

\item En el capítulo \ref{chap:herramientas}, Herramientas utilizadas, se realiza una descripción de las diferentes elementos hardware y software empleados durante el desarrollo del proyecto y necesarios para la utilización del mismo.

\item En el capítulo \ref{chap:organización-temporal}, Organización temporal, se recoge todo lo que concierne a la distribución y duración de cada una de las tareas llevadas a cabo durante el desarrollo del proyecto que el presente documento describe.

\item En el capítulo \ref{chap:dispositivos-hardware}, Configuración y montaje de los dispositivos hardware, se explica el proceso seguido para la correcta integración de los dispositivos hardware empleados describiendo la interconexión entre ellos así como su configuración. 

\item En el capítulo \ref{chap:desarrollo-software}, Desarrollo software, se realiza un análisis sobre la metodología empleada para el desarrollo software, describiendo los modelos de ciclo de vida utilizados, la descripción de los requisitos funcionales junto con el diagrama de casos de uso.

\item En el capítulo \ref{chap:reconocimiento}, Software de reconocimiento, se hace una descripción explicando los diferentes aspectos y elementos de cada uno de los prototipos desarrollados junto con los problemas encontrados y soluciones adoptadas.

\item En el capítulo \ref{chap:software-de-control}, Software de control, se describe cómo se ha llevado a cabo la comunicación ordenador-vehículo a nivel software.

\item En el capítulo \ref{chap:interfaz-gráfica}, Interfaz gráfica, se recogen aquellos aspectos técnicos de interés referentes a la elaboración de la interfaz gráfica.

\item En el capítulo \ref{chap:manual-usuario}, Guía de usuario, se describen los diferentes aspectos necesarios para la correcta utilización del conjunto software y hardware de los que se compone el presente proyecto.

\item En el capítulo \ref{chap:conclusiones}, Conclusiones, se hace mención de las conclusiones obtenidas tras la realización del proyecto además de las posibles mejoras aplicables.

\item En el capítulo Anexos \ref{chap:anexos}, aparecen los manuales de instalación del software que ha sido necesario para la realización del proyecto.

\end{itemize}
